\subsection{Įrankių ekosistema}
\label{subsec:irankiu-ekosistema}

Emuliatoriaus realizacijai pasirinkta Zig kalba. Šis pasirinkimas motyvuojamas tuo, kad emuliatorius yra žemo lygio sistema: reikia valdyti atminties išdėstymą, sveikųjų skaičių bitų pločius, MMIO prieigas, klaidų kelią ir sąsają su C bibliotekomis. Zig dokumentacijoje kalba apibūdinama kaip bendros paskirties kalba ir įrankių grandinė, skirta patikimai, optimaliai ir pakartotinai panaudojamai programinei įrangai kurti \cite{zig_language_reference}. Darbui ypač svarbus \code{comptime} mechanizmas, nes Zig leidžia dalį reikšmių ir kontrolės srautų įvertinti kompiliavimo metu \cite{zig_language_reference}. Tai naudinga generuojant tipizuotas struktūras pagal procesoriaus žodžio ilgį arba konfigūracinius plėtinius.

Kitas Zig privalumas šiame darbe yra sąsaja su C ekosistema. QPU realizacijos dalis remiasi C biblioteka, todėl emuliatorius turi patikimai susieti Zig ir C tipus. Zig dokumentacijoje aprašomas C ABI suderinamų tipų rinkinys ir C kodo vertimo galimybės \cite{zig_language_reference}. Tai leidžia emuliatoriaus branduolį rašyti Zig kalba, o specializuotas bibliotekas naudoti per aiškiai apibrėžtą tarpkalbinę sąsają.

Nix naudojamas todėl, kad darbe yra daug tarpusavyje priklausomų komponentų: emuliatorius, Zig kompiliatorius, libelf, libquantum, OpenSBI, Linux branduolys, initramfs, BusyBox, musl ir RISC-V testavimo įrankiai. Nix dokumentacijoje pabrėžiamos atkuriamos kūrimo aplinkos, deklaratyvus sistemų aprašymas, integracinis testavimas ir priklausomybių konfliktų mažinimas \cite{nix_dev}. Flake struktūra papildomai leidžia užrakinti įvestis per \code{flake.lock}; Nix dokumentacijoje \code{flake.lock} aprašomas kaip failas, kuriame laikomi \code{flake.nix} nurodytų įvesčių užraktai \cite{nix_flake_lock}. Ši savybė svarbi, nes Linux branduolio, OpenSBI ir userspace komponentų versijų nesuderinamumas gali sukelti klaidas, kurios nėra susijusios su pačiu emuliatoriumi.

Derinimui naudojama GDB nuotolinio protokolo idėja. GDB Remote Serial Protocol apibrėžia paketais pagrįstą sąsają tarp derintuvės ir nuotolinio taikinio \cite{gdb_remote_protocol}. Emuliatoriuje toks taikinys yra ne fizinis procesorius, o programinis CPU modelis. Todėl GDB sąsaja leidžia sustabdyti vykdymą, skaityti registrus, vykdyti po vieną instrukciją ir tikrinti atmintį. Tai ypač naudinga Linux paleidimo metu, kai klaida gali būti kelių sluoksnių sąveikos pasekmė: neteisingo CSR, blogo adresų vertimo, neteisingai aprašyto įrenginio arba netikslaus trap apdorojimo.
